\chapter{Технологическая часть}

В этом разделе будет описан процесс разработки программного обеспечения: используемые средства,
последовательность взаимодействия с большими языковыми моделями (БЯМ), исходные промпты,
полученные результаты, внесённые правки и их итоговый эффект.

\section{Средства реализации}

Для реализации использовался язык программиров*ания Python 3.10.
Основные библиотеки:
\begin{itemize}
  \item PyPDF2 — извлечение текста из PDF;
  \item re — регулярные выражения;
  \item csv — сохранение результатов в формате CSV;
  \item pathlib — работа с путями и каталогами;
  \item argparse — интерфейс командной строки.
\end{itemize}

\subsection{Организация итоговой программы}

Во всех финальных версиях программ выделены две части:
\begin{itemize}
  \item \textbf{алгоритм (ядро)} --- извлечение строк, выделение оглавления, проверка правил варианта,
  вычисление координат первого нарушения;
  \item \textbf{вывод и CLI} --- чтение аргументов, пакетная обработка файлов и экспорт CSV.
\end{itemize}

Такое разделение использовано для того, чтобы алгоритм можно было переиспользовать
как функцию в других программах, а формат вывода (CSV или печать в консоль) изменять независимо.

\section{Взаимодействие с БЯМ и итеративная доработка}

Разработка велась в итеративном режиме. На каждом шаге формулировался промпт с требованиями варианта,
полученный код проверялся на приложенных PDF-файлах, после чего фиксировались расхождения
и формулировались уточняющие запросы либо выполнялись ручные правки.

Ниже приведены основные попытки для каждой из трёх БЯМ.

\subsection{Попытки с ChatGPT (GPT-5.2)}

\textbf{Промпт 1.}
\begin{quote}
Сгенерируй Python-программу с использованием PyPDF2 и регулярных выражений.
Вариант: проверить отсутствие структурных элементов ВВЕДЕНИЕ, ЗАКЛЮЧЕНИЕ,
СПИСОК ИСПОЛЬЗОВАННЫХ ИСТОЧНИКОВ в тексте документа и их присутствие в содержании,
а также отсутствие ошибочных наименований. Функция должна возвращать (bool, список кортежей
(строка, (страница, строка))).
\end{quote}

\textbf{Результат.} Была получена базовая реализация с постраничным извлечением текста и поиском совпадений.
Функция возвращала булев результат и список совпадений.

\textbf{Проблема.} При запуске на тестовых PDF выявилась необходимость:
\begin{itemize}
  \item учитывать «растянутые» заголовки (пример: \texttt{В В Е Д Е Н И Е});
  \item корректно выделять область оглавления, а не искать элементы по всему документу;
  \item фиксировать координату \emph{первого нарушения} (а не только успешные совпадения).
\end{itemize}

\textbf{Правка 1.} Добавлены регулярные выражения с допуском пробелов между буквами
(через шаблон вида \texttt{V\textbackslash s*V\textbackslash s*E...}),
а также реализовано выделение области оглавления по заголовку
\texttt{СОДЕРЖАНИЕ/ОГЛАВЛЕНИЕ} и форме строк-элементов оглавления.

\textbf{Результат правки 1.} Улучшилась устойчивость к артефактам извлечения текста,
а также появилась корректная проверка «в оглавлении / в тексте».

\textbf{Правка 2.} Реализован режим пакетной обработки каталога и экспорт результатов в CSV.
Вывод унифицирован под требование: \textbf{«найдена ошибка»} / \textbf{«ошибка не найдена»}.

\textbf{Результат правки 2.} Программа стала формировать таблицу результатов автоматически,
что позволяет напрямую использовать CSV в исследовательском разделе отчёта.

\subsection{Попытки с DeepSeek DeepThink}

\textbf{Промпт 1.}
\begin{quote}
Сгенерируй альтернативную реализацию того же варианта проверки структуры PDF.
Нужны регулярные выражения, выделение оглавления, координаты (страница, строка),
экспорт CSV. Код должен отличаться по стилю от предыдущего решения.
\end{quote}

\textbf{Результат.} DeepSeek предложил более компактный код,
в котором регулярные выражения формировались как строки-паттерны,
а функции проверки были сведены к небольшому набору предикатов.

\textbf{Проблема.} В первой версии отсутствовала пакетная обработка “каталог/набор файлов”,
а также не был унифицирован формат статуса в CSV.

\textbf{Правка 1.} Добавлен отбор входных PDF по “обязательному набору” (если найден),
иначе — обработка всех \texttt{*.pdf} в каталоге.
Формат статуса приведён к точным строкам:
\textbf{«найдена ошибка»} / \textbf{«ошибка не найдена»}.

\textbf{Результат правки 1.} Реализация стала совместимой по формату вывода с другими версиями,
что необходимо для сравнения результатов.

\subsection{Попытки с Qwen Coder}

\textbf{Промпт 1.}
\begin{quote}
Сгенерируй реализацию варианта проверки структуры PDF в инженерном стиле:
вынеси алгоритм в отдельный класс, обеспечь чистый интерфейс, экспорт CSV и обработку каталога.
Результат и координаты как в задании.
\end{quote}

\textbf{Результат.} Qwen сформировал структуру в виде класса \texttt{VariantChecker},
где регулярные выражения и операции выделения оглавления инкапсулированы в методах.
CLI реализован минималистично.

\textbf{Проблема.} В первой версии требовалась унификация критериев “первого нарушения”:
в некоторых случаях координаты не определялись одинаково по сравнению с другими реализациями.

\textbf{Правка 1.} Приведена логика вычисления координат к общему правилу:
если проверка провалена и координату однозначно определить нельзя, используется символ «-».
Это обеспечивает корректное сравнение таблиц результатов.

\textbf{Результат правки 1.} Таблицы сравнения стали совпадать по статусам и координатам
на тестовом наборе файлов.

\section{Пример запуска и получаемые файлы}

Проверка одного файла:
\begin{verbatim}
python pdf_checker.py path/to/file.pdf --out results_chatgpt.csv
\end{verbatim}

Проверка каталога (если найден обязательный набор, берутся только он; иначе все *.pdf):
\begin{verbatim}
python pdf_checker.py . --out results_chatgpt.csv
python pdf_checker1_deepseek.py . --out results_deepseek.csv
python pdf_checker2_qwen.py . --out results_qwen.csv
\end{verbatim}

Формируемая таблица CSV содержит колонки: \texttt{file}, \texttt{status}, \texttt{first\_coord},
где \texttt{status} принимает значения \textbf{«найдена ошибка»} или \textbf{«ошибка не найдена»}.

\section{Вывод}

В этом разделе было описано применяемое программное окружение и структура итоговых программ,
а также приведён итеративный процесс разработки с использованием трёх БЯМ:
исходные промпты, полученные результаты, выявленные несоответствия и внесённые правки,
которые привели к финальным реализациям с единым форматом вывода и возможностью сравнения результатов.
